% --- Archon physics formalization source begin ---
% archon:physics
% archon:covers PhyXMiniProblems/problem_phyx_mini_0637.lean
% archon:source-report reports/phyx_mini/problem_phyx_mini_0637.source.json

\chapter{Physics problem phyx\_mini\_0637}
\label{ch:PhyXMiniProblems_problem_phyx_mini_0637}

\paragraph{Problem source.}
\#\# Physical scenario

Fluorescence is the absorption of light at one wavelength followed by emission at a longer wavelength. Suppose a hydrogen atom in its ground state absorbs an ultraviolet photon with a wavelength of \$95.10\textbackslash{};nm\$. Immediately after the absorption, the atom undergoes a quantum jump with \$\textbackslash{}Delta n = 3\$.

\#\# Question

What is the wavelength of the photon emitted in this quantum jump?

\#\# Answer choices

- A: A: 424 nm

- B: B: 434 nm

- C: C: 444 nm

- D: D: 454 nm

\#\# Figure caption (auxiliary; use the image as primary evidence)

The image depicts an energy level diagram often used in quantum mechanics to illustrate electron transitions in an atom. 

- **Horizontal Lines**: Multiple horizontal lines represent different energy levels, labeled with "n" and the "Ground state" at the bottom. The top is labeled "Ionization limit" with a dashed line.
  
- **Vertical Arrows**: 
  - An upward arrow labeled "Absorption" indicates the process where an electron absorbs energy and moves to a higher energy level.
  - A downward arrow labeled "Emission" shows the process where an electron releases energy and returns to a lower energy level.

- **Wavy Arrows**: Associated with both absorption and emission, these wavy lines represent the energy (or photons) involved in the transitions.

Overall, this diagram explains the basic principles of electron excitation and relaxation in atomic physics.

\#\# Dataset metadata

Quantum Phenomena; Physical Model Grounding Reasoning, Multi-Formula Reasoning

\paragraph{Recorded answer/context.}
B: B: 434 nm

\paragraph{Figure/image path.}
/root/proposal\_for\_physic/science-mango/phyx\_mini\_run/phyx\_data/test\_image/637.png

\paragraph{Formalization target.}
create a compiling Lean file with sorry bodies at `PhyXMiniProblems/problem\_phyx\_mini\_0637.lean`. The Lean declarations must preserve the physical quantities, dimensions or dimensional roles, figure labels, governing-law hypotheses, and final relation expressed by this problem.
Use Mathlib/Physlib names found through LeanExplore where available. If a physics API is missing, introduce faithful local abstractions rather than scalar placeholder aliases.

\begin{theorem}[Physics formalization target]
\label{thm:physics:phyx_mini_0637:target}
\lean{PhyXMiniProblems.ProblemPhyXMini0637.emittedPhotonWavelength_is_recordedChoiceB}
\uses{def:physics:phyx-mini-0637:phyxminiproblems-problemphyxmini0637-matchessuppliedfluorescencefigure, def:physics:phyx-mini-0637:phyxminiproblems-problemphyxmini0637-hydrogenfluorescencesetup, def:physics:phyx-mini-0637:phyxminiproblems-problemphyxmini0637-matcheshydrogenfluorescencescenario, def:physics:phyx-mini-0637:phyxminiproblems-problemphyxmini0637-matchesabsorbedphotonmeasurement, def:physics:phyx-mini-0637:phyxminiproblems-problemphyxmini0637-usesstandardhydrogenreferencedata, def:physics:phyx-mini-0637:phyxminiproblems-problemphyxmini0637-hasphysicalhydrogenfluorescenceparameters, def:physics:phyx-mini-0637:phyxminiproblems-problemphyxmini0637-satisfieshydrogenfluorescencelaws, def:physics:phyx-mini-0637:phyxminiproblems-problemphyxmini0637-usesroundedabsorptionlevelidentification, def:physics:phyx-mini-0637:phyxminiproblems-problemphyxmini0637-recordeddatasetanswer, def:physics:phyx-mini-0637:phyxminiproblems-problemphyxmini0637-roundstodisplayedwavelength, def:physics:phyx-mini-0637:phyxminiproblems-problemphyxmini0637-isuniqueclosestdisplayedwavelength}
The assigned autoformalize agent should translate this physics problem into Lean declarations in the covered file, with theorem and lemma proof bodies written as `by sorry`.
\end{theorem}
\begin{proof}
This is an autoformalization task, not a proof task. Produce statements that can later be proved by the physics prover without weakening the physical model.
\end{proof}

% --- Archon declaration topology begin ---
\section{Declaration topology}
\begin{definition}[Wavelength Quantity]
  \label{def:physics:phyx-mini-0637:phyxminiproblems-problemphyxmini0637-wavelengthquantity}
  \lean{PhyXMiniProblems.ProblemPhyXMini0637.WavelengthQuantity}
  A nonnegative, unit-independent physical length used as a photon wavelength
\end{definition}
\begin{proof}
  This is the declared model definition of Wavelength Quantity.
\end{proof}
\begin{definition}[energy In Joules]
  \label{def:physics:phyx-mini-0637:phyxminiproblems-problemphyxmini0637-energyinjoules}
  \lean{PhyXMiniProblems.ProblemPhyXMini0637.energyInJoules}
  Read a unit-independent energy in coherent-SI joules
\end{definition}
\begin{proof}
  This is the declared model definition of energy In Joules.
\end{proof}
\begin{definition}[energy In Electron Volts]
  \label{def:physics:phyx-mini-0637:phyxminiproblems-problemphyxmini0637-energyinelectronvolts}
  \lean{PhyXMiniProblems.ProblemPhyXMini0637.energyInElectronVolts}
  \uses{def:physics:phyx-mini-0637:phyxminiproblems-problemphyxmini0637-energyinjoules}
  Read an energy in electron volts using Physlib's calibrated electron volt
\end{definition}
\begin{proof}
  This is the declared model definition of energy In Electron Volts.
\end{proof}
\begin{definition}[wavelength Readout]
  \label{def:physics:phyx-mini-0637:phyxminiproblems-problemphyxmini0637-wavelengthreadout}
  \lean{PhyXMiniProblems.ProblemPhyXMini0637.wavelengthReadout}
  \uses{def:physics:phyx-mini-0637:phyxminiproblems-problemphyxmini0637-wavelengthquantity}
  Read a physical wavelength in a selected unit of length
\end{definition}
\begin{proof}
  This is the declared model definition of wavelength Readout.
\end{proof}
\begin{definition}[wavelength In Meters]
  \label{def:physics:phyx-mini-0637:phyxminiproblems-problemphyxmini0637-wavelengthinmeters}
  \lean{PhyXMiniProblems.ProblemPhyXMini0637.wavelengthInMeters}
  \uses{def:physics:phyx-mini-0637:phyxminiproblems-problemphyxmini0637-wavelengthquantity, def:physics:phyx-mini-0637:phyxminiproblems-problemphyxmini0637-wavelengthreadout}
  Metre readout of a photon wavelength
\end{definition}
\begin{proof}
  This is the declared model definition of wavelength In Meters.
\end{proof}
\begin{definition}[wavelength In Nanometers]
  \label{def:physics:phyx-mini-0637:phyxminiproblems-problemphyxmini0637-wavelengthinnanometers}
  \lean{PhyXMiniProblems.ProblemPhyXMini0637.wavelengthInNanometers}
  \uses{def:physics:phyx-mini-0637:phyxminiproblems-problemphyxmini0637-wavelengthquantity, def:physics:phyx-mini-0637:phyxminiproblems-problemphyxmini0637-wavelengthreadout}
  Nanometre readout used for the measured and displayed wavelengths
\end{definition}
\begin{proof}
  This is the declared model definition of wavelength In Nanometers.
\end{proof}
\begin{definition}[planck Times Light In Electron Volt Nanometers]
  \label{def:physics:phyx-mini-0637:phyxminiproblems-problemphyxmini0637-plancktimeslightinelectronvoltnanometers}
  \lean{PhyXMiniProblems.ProblemPhyXMini0637.planckTimesLightInElectronVoltNanometers}
  \uses{def:physics:phyx-mini-0637:phyxminiproblems-problemphyxmini0637-energyinjoules}
  The product h c in electron-volt nanometres. Physlib supplies the reduced Planck constant ℏ, the SI speed of light, and the electron-volt calibration; 2π converts ℏ to h, while 10\textasciicircum{}9 converts metres to nanometres
\end{definition}
\begin{proof}
  This is the declared model definition of planck Times Light In Electron Volt Nanometers.
\end{proof}
\begin{definition}[Figure Text Label]
  \label{def:physics:phyx-mini-0637:phyxminiproblems-problemphyxmini0637-figuretextlabel}
  \lean{PhyXMiniProblems.ProblemPhyXMini0637.FigureTextLabel}
  Text labels visibly printed in the raster
\end{definition}
\begin{proof}
  This is the declared model definition of Figure Text Label.
\end{proof}
\begin{definition}[Transition Arrow]
  \label{def:physics:phyx-mini-0637:phyxminiproblems-problemphyxmini0637-transitionarrow}
  \lean{PhyXMiniProblems.ProblemPhyXMini0637.TransitionArrow}
  The two straight transition arrows in the diagram
\end{definition}
\begin{proof}
  This is the declared model definition of Transition Arrow.
\end{proof}
\begin{definition}[Hydrogen Fluorescence Figure]
  \label{def:physics:phyx-mini-0637:phyxminiproblems-problemphyxmini0637-hydrogenfluorescencefigure}
  \lean{PhyXMiniProblems.ProblemPhyXMini0637.HydrogenFluorescenceFigure}
  \uses{def:physics:phyx-mini-0637:phyxminiproblems-problemphyxmini0637-figuretextlabel, def:physics:phyx-mini-0637:phyxminiproblems-problemphyxmini0637-transitionarrow}
  Presentation data extracted from image 637.png. It deliberately records no numerical emitted wavelength and does not identify either endpoint of the emission arrow with a particular principal quantum number
\end{definition}
\begin{proof}
  This is the declared model definition of Hydrogen Fluorescence Figure.
\end{proof}
\begin{definition}[Matches Supplied Fluorescence Figure]
  \label{def:physics:phyx-mini-0637:phyxminiproblems-problemphyxmini0637-matchessuppliedfluorescencefigure}
  \lean{PhyXMiniProblems.ProblemPhyXMini0637.MatchesSuppliedFluorescenceFigure}
  \uses{def:physics:phyx-mini-0637:phyxminiproblems-problemphyxmini0637-hydrogenfluorescencefigure}
  Literal readouts from the primary raster. The image shows six solid upper level lines in addition to the ground line, a dashed ionization limit, an upward absorption arrow, and a downward emission arrow
\end{definition}
\begin{proof}
  This is the declared model definition of Matches Supplied Fluorescence Figure.
\end{proof}
\begin{definition}[Hydrogen Fluorescence Setup]
  \label{def:physics:phyx-mini-0637:phyxminiproblems-problemphyxmini0637-hydrogenfluorescencesetup}
  \lean{PhyXMiniProblems.ProblemPhyXMini0637.HydrogenFluorescenceSetup}
  \uses{def:physics:phyx-mini-0637:phyxminiproblems-problemphyxmini0637-wavelengthquantity, def:physics:phyx-mini-0637:phyxminiproblems-problemphyxmini0637-hydrogenfluorescencefigure}
  The physical observables and state labels for the experiment. In particular, the emitted wavelength is an independent dimensionful field; it is not defined from an answer choice or from the requested value
\end{definition}
\begin{proof}
  This is the declared model definition of Hydrogen Fluorescence Setup.
\end{proof}
\begin{definition}[Matches Hydrogen Fluorescence Scenario]
  \label{def:physics:phyx-mini-0637:phyxminiproblems-problemphyxmini0637-matcheshydrogenfluorescencescenario}
  \lean{PhyXMiniProblems.ProblemPhyXMini0637.MatchesHydrogenFluorescenceScenario}
  \uses{def:physics:phyx-mini-0637:phyxminiproblems-problemphyxmini0637-wavelengthinnanometers, def:physics:phyx-mini-0637:phyxminiproblems-problemphyxmini0637-hydrogenfluorescencesetup}
  Qualitative facts stated in the problem, including the downward Δn = 3 jump
\end{definition}
\begin{proof}
  This is the declared model definition of Matches Hydrogen Fluorescence Scenario.
\end{proof}
\begin{definition}[Matches Absorbed Photon Measurement]
  \label{def:physics:phyx-mini-0637:phyxminiproblems-problemphyxmini0637-matchesabsorbedphotonmeasurement}
  \lean{PhyXMiniProblems.ProblemPhyXMini0637.MatchesAbsorbedPhotonMeasurement}
  \uses{def:physics:phyx-mini-0637:phyxminiproblems-problemphyxmini0637-wavelengthinnanometers, def:physics:phyx-mini-0637:phyxminiproblems-problemphyxmini0637-hydrogenfluorescencesetup}
  The ultraviolet wavelength reported in the prose, as a nanometre readout
\end{definition}
\begin{proof}
  This is the declared model definition of Matches Absorbed Photon Measurement.
\end{proof}
\begin{definition}[Uses Standard Hydrogen Reference Data]
  \label{def:physics:phyx-mini-0637:phyxminiproblems-problemphyxmini0637-usesstandardhydrogenreferencedata}
  \lean{PhyXMiniProblems.ProblemPhyXMini0637.UsesStandardHydrogenReferenceData}
  \uses{def:physics:phyx-mini-0637:phyxminiproblems-problemphyxmini0637-energyinelectronvolts, def:physics:phyx-mini-0637:phyxminiproblems-problemphyxmini0637-hydrogenfluorescencesetup}
  Standard ground-level calibration used by the elementary Bohr spectrum. This states E₁ = -13.6 eV; it does not state the excited level, emission endpoint, or emitted wavelength
\end{definition}
\begin{proof}
  This is the declared model definition of Uses Standard Hydrogen Reference Data.
\end{proof}
\begin{definition}[Has Physical Hydrogen Fluorescence Parameters]
  \label{def:physics:phyx-mini-0637:phyxminiproblems-problemphyxmini0637-hasphysicalhydrogenfluorescenceparameters}
  \lean{PhyXMiniProblems.ProblemPhyXMini0637.HasPhysicalHydrogenFluorescenceParameters}
  \uses{def:physics:phyx-mini-0637:phyxminiproblems-problemphyxmini0637-energyinelectronvolts, def:physics:phyx-mini-0637:phyxminiproblems-problemphyxmini0637-wavelengthinmeters, def:physics:phyx-mini-0637:phyxminiproblems-problemphyxmini0637-hydrogenfluorescencesetup}
  Positivity and bound-state conditions for the physical branch of the model
\end{definition}
\begin{proof}
  This is the declared model definition of Has Physical Hydrogen Fluorescence Parameters.
\end{proof}
\begin{definition}[Satisfies Hydrogen Fluorescence Laws]
  \label{def:physics:phyx-mini-0637:phyxminiproblems-problemphyxmini0637-satisfieshydrogenfluorescencelaws}
  \lean{PhyXMiniProblems.ProblemPhyXMini0637.SatisfiesHydrogenFluorescenceLaws}
  \uses{def:physics:phyx-mini-0637:phyxminiproblems-problemphyxmini0637-energyinelectronvolts, def:physics:phyx-mini-0637:phyxminiproblems-problemphyxmini0637-wavelengthinnanometers, def:physics:phyx-mini-0637:phyxminiproblems-problemphyxmini0637-plancktimeslightinelectronvoltnanometers, def:physics:phyx-mini-0637:phyxminiproblems-problemphyxmini0637-hydrogenfluorescencesetup}
  The elementary hydrogen spectrum, Planck--Einstein photon relations, and emission energy balance. Each law is uniform or relates independent physical quantities; no field contains n = 5, n = 2, 434 nm, or an answer label
\end{definition}
\begin{proof}
  This is the declared model definition of Satisfies Hydrogen Fluorescence Laws.
\end{proof}
\begin{definition}[post Absorption Energy Estimate Electron Volts]
  \label{def:physics:phyx-mini-0637:phyxminiproblems-problemphyxmini0637-postabsorptionenergyestimateelectronvolts}
  \lean{PhyXMiniProblems.ProblemPhyXMini0637.postAbsorptionEnergyEstimateElectronVolts}
  \uses{def:physics:phyx-mini-0637:phyxminiproblems-problemphyxmini0637-energyinelectronvolts, def:physics:phyx-mini-0637:phyxminiproblems-problemphyxmini0637-hydrogenfluorescencesetup}
  Energy obtained by adding the absorbed photon energy to the initial atomic energy. It is a scalar electron-volt readout used only to interpret the rounded wavelength datum
\end{definition}
\begin{proof}
  This is the declared model definition of post Absorption Energy Estimate Electron Volts.
\end{proof}
\begin{definition}[Is Unique Closest Post Absorption Level]
  \label{def:physics:phyx-mini-0637:phyxminiproblems-problemphyxmini0637-isuniqueclosestpostabsorptionlevel}
  \lean{PhyXMiniProblems.ProblemPhyXMini0637.IsUniqueClosestPostAbsorptionLevel}
  \uses{def:physics:phyx-mini-0637:phyxminiproblems-problemphyxmini0637-energyinelectronvolts, def:physics:phyx-mini-0637:phyxminiproblems-problemphyxmini0637-hydrogenfluorescencesetup, def:physics:phyx-mini-0637:phyxminiproblems-problemphyxmini0637-postabsorptionenergyestimateelectronvolts}
  A positive principal level is the unique closest discrete hydrogen level to the energy estimate inferred from the rounded 95.10 nm datum
\end{definition}
\begin{proof}
  This is the declared model definition of Is Unique Closest Post Absorption Level.
\end{proof}
\begin{definition}[Uses Rounded Absorption Level Identification]
  \label{def:physics:phyx-mini-0637:phyxminiproblems-problemphyxmini0637-usesroundedabsorptionlevelidentification}
  \lean{PhyXMiniProblems.ProblemPhyXMini0637.UsesRoundedAbsorptionLevelIdentification}
  \uses{def:physics:phyx-mini-0637:phyxminiproblems-problemphyxmini0637-hydrogenfluorescencesetup, def:physics:phyx-mini-0637:phyxminiproblems-problemphyxmini0637-isuniqueclosestpostabsorptionlevel}
  Interpretation rule for the rounded absorption measurement. It selects the closest positive stationary level without naming which level that is
\end{definition}
\begin{proof}
  This is the declared model definition of Uses Rounded Absorption Level Identification.
\end{proof}
\begin{lemma}[absorbed Photon identifies excited Level Five]
  \label{lem:physics:phyx-mini-0637:phyxminiproblems-problemphyxmini0637-absorbedphoton-identifies-excitedlevelfive}
  \lean{PhyXMiniProblems.ProblemPhyXMini0637.absorbedPhoton_identifies_excitedLevelFive}
  \uses{def:physics:phyx-mini-0637:phyxminiproblems-problemphyxmini0637-hydrogenfluorescencesetup, def:physics:phyx-mini-0637:phyxminiproblems-problemphyxmini0637-matcheshydrogenfluorescencescenario, def:physics:phyx-mini-0637:phyxminiproblems-problemphyxmini0637-matchesabsorbedphotonmeasurement, def:physics:phyx-mini-0637:phyxminiproblems-problemphyxmini0637-usesstandardhydrogenreferencedata, def:physics:phyx-mini-0637:phyxminiproblems-problemphyxmini0637-hasphysicalhydrogenfluorescenceparameters, def:physics:phyx-mini-0637:phyxminiproblems-problemphyxmini0637-satisfieshydrogenfluorescencelaws, def:physics:phyx-mini-0637:phyxminiproblems-problemphyxmini0637-usesroundedabsorptionlevelidentification}
  The rounded absorption datum and standard spectrum identify the excited level
\end{lemma}
\begin{proof}
  The conclusion follows from the governing relations and figure readouts named in the statement of absorbed Photon identifies excited Level Five.
\end{proof}
\begin{lemma}[emission Transition is five To Two]
  \label{lem:physics:phyx-mini-0637:phyxminiproblems-problemphyxmini0637-emissiontransition-is-fivetotwo}
  \lean{PhyXMiniProblems.ProblemPhyXMini0637.emissionTransition_is_fiveToTwo}
  \uses{def:physics:phyx-mini-0637:phyxminiproblems-problemphyxmini0637-hydrogenfluorescencesetup, def:physics:phyx-mini-0637:phyxminiproblems-problemphyxmini0637-matcheshydrogenfluorescencescenario, def:physics:phyx-mini-0637:phyxminiproblems-problemphyxmini0637-matchesabsorbedphotonmeasurement, def:physics:phyx-mini-0637:phyxminiproblems-problemphyxmini0637-usesstandardhydrogenreferencedata, def:physics:phyx-mini-0637:phyxminiproblems-problemphyxmini0637-hasphysicalhydrogenfluorescenceparameters, def:physics:phyx-mini-0637:phyxminiproblems-problemphyxmini0637-satisfieshydrogenfluorescencelaws, def:physics:phyx-mini-0637:phyxminiproblems-problemphyxmini0637-usesroundedabsorptionlevelidentification}
  Combining the identified excited state with Δn = 3 gives the 5 → 2 jump
\end{lemma}
\begin{proof}
  The conclusion follows from the governing relations and figure readouts named in the statement of emission Transition is five To Two.
\end{proof}
\begin{lemma}[emitted Photon Wavelength eq planck Times Light div level Gap]
  \label{lem:physics:phyx-mini-0637:phyxminiproblems-problemphyxmini0637-emittedphotonwavelength-eq-plancktimeslight-div-levelgap}
  \lean{PhyXMiniProblems.ProblemPhyXMini0637.emittedPhotonWavelength_eq_planckTimesLight_div_levelGap}
  \uses{def:physics:phyx-mini-0637:phyxminiproblems-problemphyxmini0637-energyinelectronvolts, def:physics:phyx-mini-0637:phyxminiproblems-problemphyxmini0637-wavelengthinnanometers, def:physics:phyx-mini-0637:phyxminiproblems-problemphyxmini0637-plancktimeslightinelectronvoltnanometers, def:physics:phyx-mini-0637:phyxminiproblems-problemphyxmini0637-hydrogenfluorescencesetup, def:physics:phyx-mini-0637:phyxminiproblems-problemphyxmini0637-hasphysicalhydrogenfluorescenceparameters, def:physics:phyx-mini-0637:phyxminiproblems-problemphyxmini0637-satisfieshydrogenfluorescencelaws}
  Solving the emission photon law expresses its wavelength through the physical level-energy gap. This helper does not assign a displayed numeric answer
\end{lemma}
\begin{proof}
  The conclusion follows from the governing relations and figure readouts named in the statement of emitted Photon Wavelength eq planck Times Light div level Gap.
\end{proof}
\begin{definition}[Answer Choice]
  \label{def:physics:phyx-mini-0637:phyxminiproblems-problemphyxmini0637-answerchoice}
  \lean{PhyXMiniProblems.ProblemPhyXMini0637.AnswerChoice}
  Labels of the four multiple-choice answers
\end{definition}
\begin{proof}
  This is the declared model definition of Answer Choice.
\end{proof}
\begin{definition}[wavelength Nanometers]
  \label{def:physics:phyx-mini-0637:phyxminiproblems-problemphyxmini0637-answerchoice-wavelengthnanometers}
  \lean{PhyXMiniProblems.ProblemPhyXMini0637.AnswerChoice.wavelengthNanometers}
  \uses{def:physics:phyx-mini-0637:phyxminiproblems-problemphyxmini0637-answerchoice}
  Wavelength printed next to each answer label, in nanometres
\end{definition}
\begin{proof}
  This is the declared model definition of wavelength Nanometers.
\end{proof}
\begin{definition}[recorded Dataset Answer]
  \label{def:physics:phyx-mini-0637:phyxminiproblems-problemphyxmini0637-recordeddatasetanswer}
  \lean{PhyXMiniProblems.ProblemPhyXMini0637.recordedDatasetAnswer}
  \uses{def:physics:phyx-mini-0637:phyxminiproblems-problemphyxmini0637-answerchoice}
  The answer label recorded by the supplied dataset metadata
\end{definition}
\begin{proof}
  This is the declared model definition of recorded Dataset Answer.
\end{proof}
\begin{definition}[Rounds To Displayed Wavelength]
  \label{def:physics:phyx-mini-0637:phyxminiproblems-problemphyxmini0637-roundstodisplayedwavelength}
  \lean{PhyXMiniProblems.ProblemPhyXMini0637.RoundsToDisplayedWavelength}
  \uses{def:physics:phyx-mini-0637:phyxminiproblems-problemphyxmini0637-wavelengthquantity, def:physics:phyx-mini-0637:phyxminiproblems-problemphyxmini0637-wavelengthinnanometers, def:physics:phyx-mini-0637:phyxminiproblems-problemphyxmini0637-answerchoice}
  A physical wavelength rounds to the displayed whole-nanometre value
\end{definition}
\begin{proof}
  This is the declared model definition of Rounds To Displayed Wavelength.
\end{proof}
\begin{definition}[Is Unique Closest Displayed Wavelength]
  \label{def:physics:phyx-mini-0637:phyxminiproblems-problemphyxmini0637-isuniqueclosestdisplayedwavelength}
  \lean{PhyXMiniProblems.ProblemPhyXMini0637.IsUniqueClosestDisplayedWavelength}
  \uses{def:physics:phyx-mini-0637:phyxminiproblems-problemphyxmini0637-wavelengthquantity, def:physics:phyx-mini-0637:phyxminiproblems-problemphyxmini0637-wavelengthinnanometers, def:physics:phyx-mini-0637:phyxminiproblems-problemphyxmini0637-answerchoice}
  A displayed answer is strictly nearer to the physical wavelength than every alternative
\end{definition}
\begin{proof}
  This is the declared model definition of Is Unique Closest Displayed Wavelength.
\end{proof}
% --- Archon declaration topology end ---
% --- Archon physics formalization source end ---
